Para todo inteiro positivo \(n\), seja \(p(n)\) o número de maneiras de expressar \(n\) como uma soma de inteiros positivos. Por exemplo, \(p(4)=5\) porque

\[4=3+1=2+2=2+1+1=1+1+1.\]

Além disso, definimos \(p(0)=1\).

Prove que \(p(n)-p(n-1)\) é o número de maneiras de expressar \(n\) como uma soma de inteiros cada um dos quais é estritamente maior que 1.
